\documentclass[10pt, twocolumn]{article}

% --- Essential Packages ---
\usepackage[margin=1.5cm]{geometry} % Narrow margins typical of journals
\usepackage{graphicx}               % For figures
\usepackage{authblk}                % For author affiliations
\usepackage{caption}                % To format figure captions
\usepackage{amsmath, amssymb, amsfonts} % For math environments
\usepackage{titlesec}               % To format section headings
\usepackage{booktabs}               % For tables
\usepackage{xcolor}                 % For TODO highlighting (load before hyperref)

% --- Font and Aesthetic Tweaks ---
\usepackage{mathptmx}               % Times-like font
\usepackage[T1]{fontenc}
\usepackage[colorlinks=true, allcolors=blue]{hyperref}

% --- Formatting Section Headings ---
\titleformat{\section}{\large\bfseries}{\thesection.}{0.5em}{}
\titleformat{\subsection}{\normalsize\bfseries}{\thesubsection.}{0.5em}{}
\titleformat{\subsubsection}{\normalsize\itshape}{\thesubsubsection.}{0.5em}{}

% --- Formatting the Abstract ---
\renewenvironment{abstract}
 {\small\bfseries}
 {\par\vspace{1em}}

% --- Helpers (DRAFT only: remove before submission) ---
\newcommand{\todo}[1]{\textcolor{red}{[\textbf{TODO}: #1]}}
% Compilable figure placeholder; replace the box with \includegraphics when ready.
\newcommand{\figplaceholder}[1]{%
  \fbox{\parbox[c][3.0cm][c]{0.92\linewidth}{\centering\small\itshape #1}}}

% --- Title and Author Information ---
\title{\vspace{-1.5cm} \textbf{MOGFormer: A Patient-Specific Multi-Omics Graph Transformer with Hierarchical Pathway Encodings for Interpretable Cancer Subtyping}}

\author[1]{BOBANT Jean-Baptiste\thanks{Corresponding author: bobantjeanbaptiste@gmail.com}}

\affil[1]{\small CentraleSupélec, Paris-Saclay University, France}
\date{}

\begin{document}

\maketitle

% =====================================================================
% ABSTRACT
% =====================================================================
\section*{Abstract}
Accurate stratification of cancer subtypes is essential for personalized
therapy, yet deep models struggle in the $N \ll p$ regime characteristic of
multi-omics cohorts, where well-tuned classical early-fusion models remain
formidable baselines. Multi-modal fusion methods also suffer from information
routing biases, lazily relying on high-density transcriptomics while ignoring
epigenetic and copy-number signals, and standard graph neural networks fail to
respect established biological topographies. We present MOGFormer, a
patient-specific multi-omics graph transformer that builds unified gene-level
tokens from raw multimodal continuous signals through non-linear feature
tokenizers and localized intra-gene cross-attention, and structures its global
attention with a soft inductive bias derived from Protein--Protein Interaction
(PPI) shortest-path distances, read out through specialized Latent Pathway
Tokens. We evaluate MOGFormer on $949$ TCGA-BRCA patients (RNA-seq, copy-number variation,
methylation) for PAM50 subtyping under a deliberately rigorous, leakage-free
protocol: repeated stratified patient-level cross-validation with all feature
selection, normalization, and graph construction refit inside each fold, and
paired comparison against a full classical baseline suite. Beyond predictive
performance, MOGFormer offers native, multi-tiered interpretability, mapping
global patient states back to single-gene multimodal driver mechanisms without
post-hoc explainers.

% =====================================================================
% INTRODUCTION
% =====================================================================
\section*{Introduction}

In the era of precision oncology, leveraging machine learning for \textit{in
silico} biological discovery requires architectures that simultaneously navigate
vast feature spaces and yield mechanistic insight. Traditional early-fusion
models often match or outperform deep networks on multi-omics cohorts because of
the acute $N \ll p$ regime, where the absence of structural constraints lets
flexible models overfit~\cite{gorishniy2022,wang2021mogonet}. To embed biological
reality into deep models, Graph Neural Networks (GNNs) have been used with Gene
Regulatory Networks or PPI networks as structural priors. However, message-passing
GNNs suffer from two architectural flaws: the over-squashing of information during
localized aggregation, and an inability to discover \textit{de novo} or
context-specific interactions beyond the static edges of the input graph.

To transcend these limitations, we extend the transformer architecture to
biological graphs through \texttt{MOGFormer}. Following the Graphormer
paradigm~\cite{ying2021graphormer}, global self-attention parametrized by a soft
structural bias prevents information bottlenecks while letting the network learn
long-range genetic dependencies absent from static databases.

Crucially, cellular state is not driven by transcriptomics alone; it is governed
by a multi-tiered epigenetic and structural landscape. \texttt{MOGFormer}
captures three primary modalities:
\begin{itemize}
    \item \textbf{RNA-sequencing (mRNA):} the high-density downstream
    transcriptional output and functional state of the cell.
    \item \textbf{Copy Number Variations (CNV):} discrete structural alterations
    (focal amplifications, deletions) that set baseline expression boundaries.
    \item \textbf{Methylation $\beta$-values:} localized CpG-island modifications
    that dynamically silence or activate transcription.
\end{itemize}
Rather than treating these as disjoint vectors, our architecture explicitly models
their intra-gene cross-talk (e.g.\ methylation-driven silencing of expression)
prior to global state computation.

We validate this framework on the The Cancer Genome Atlas Breast Invasive
Carcinoma (TCGA-BRCA) cohort~\cite{tcga2012brca}. Breast cancer is highly
heterogeneous, conventionally classified via the PAM50
panel~\cite{parker2009pam50} into Basal-like, HER2-enriched, Luminal A, Luminal B,
and Normal-like subtypes, which carry distinct prognosis and therapeutic
implications. While robust classification matters, the core contribution of
\texttt{MOGFormer} lies in its native, multi-tiered interpretability: bypassing
post-hoc explainers (LRP, SHAP), the routing mechanics expose which global
pathways drove a classification, which driver genes anchored those pathways, and
which omics modality dictated an anomalous gene state.

% =====================================================================
% DATA
% =====================================================================
\section{Dataset and Preprocessing}

\subsection{Dataset}
We use TCGA-BRCA, retaining female patients with matched RNA-seq, copy-number, and
methylation $\beta$-value profiles for the genes of interest, yielding $949$
patients. The PAM50 label distribution is markedly imbalanced: 166 Basal, 77 HER2,
480 Luminal A, 194 Luminal B, and 35 Normal-like (Fig.~\ref{fig:classdist}). This
imbalance, in particular the 35 Normal-like cases, governs both the
evaluation design (Sec.~\ref{sec:eval}) and the dominant error modes
(Sec.~\ref{sec:results}).

\begin{figure}[t]
    \centering
    % \includegraphics[width=0.95\linewidth]{figures/V01_class_distribution.pdf}
    \figplaceholder{Figure 1 placeholder --- PAM50 class distribution (the imbalance), 166/77/480/194/35.}
    \caption{PAM50 subtype distribution across the 949-patient TCGA-BRCA cohort.}
    \label{fig:classdist}
\end{figure}

\subsection{Leakage-Free, Per-Fold Preprocessing}
Preventing data leakage is critical and is enforced \emph{within} every
cross-validation fold rather than once globally. Let the raw inputs be
$\mathbf{X}_{mRNA}$, $\mathbf{X}_{CNV}$, $\mathbf{X}_{methy}$ and the clinical
vector $\mathbf{Y}$. Only patients present across all three modalities and the
clinical table are retained, deduplicated at the patient level. Every statistical
operation below is fit on the training fold alone and merely applied to held-out
patients.

\vspace{0.2cm}\noindent\textbf{Highly Variable Gene (HVG) Selection.}
To extract informative genes while remaining robust to outliers, we rank by the
Median Absolute Deviation (MAD),
\begin{equation}
    \text{MAD}(\mathbf{x}) = \text{median}(|\mathbf{x} - \text{median}(\mathbf{x})|),
\end{equation}
computed per modality, and combine modalities through a consensus rank
\begin{equation}
    \text{Rank}_{cons} = \text{Rank}_{mRNA} + \text{Rank}_{CNV} + \text{Rank}_{methy}.
\end{equation}
The top $N$ genes by consensus rank are kept. A curated set of PAM50/Prosigna and
known BRCA driver genes is \emph{force-included} after ranking, guaranteeing that
canonical markers (e.g.\ \textit{ESR1}, \textit{ERBB2}, \textit{MKI67}) are present
in every fold regardless of their MAD rank. Unless otherwise stated we use
$N = 250$ MAD genes augmented by the curated set.

\vspace{0.2cm}\noindent\textbf{Modality-Specific Normalization.}
Normalization parameters are fit on the training split only. For RNA-seq we apply
$\mathbf{x}_{mRNA}' = \text{StandardScale}(\ln(\mathbf{x}_{mRNA} + 1))$; for CNV
and methylation, $\mathbf{x}' = \text{StandardScale}(\mathbf{x})$. Median
imputation, where required, is likewise fit on the training fold.

% =====================================================================
% ARCHITECTURE
% =====================================================================
\section{Architectural Formulation}

\subsection{Non-Linear Modality Lifting}
Genomic and epigenomic variations do not scale linearly; biological mechanisms
operate on physiological thresholds. For each patient and selected gene
$i \in \{1,\dots,N\}$, scalar values are projected into a $d$-dimensional space
($d=128$) by per-modality Feature Tokenizers (small MLPs), with learnable modality
embeddings $\mathbf{E}_m, \mathbf{E}_c, \mathbf{E}_t \in \mathbb{R}^d$ encoding
biological origin:
\begin{align}
    \mathbf{z}_{m,i} &= \text{MLP}_{mRNA}(x_{m,i}) + \mathbf{E}_{m} \\
    \mathbf{z}_{c,i} &= \text{MLP}_{CNV}(x_{c,i}) + \mathbf{E}_{c} \\
    \mathbf{z}_{t,i} &= \text{MLP}_{methy}(x_{t,i}) + \mathbf{E}_{t}
\end{align}
Each MLP is \texttt{Linear} $\to$ \texttt{LayerNorm} $\to$ \texttt{GELU} $\to$
\texttt{Linear}.

\subsection{Intra-Gene Fusion and Symmetrical Masking}
To model local epigenetic--transcriptomic interactions (e.g.\ methylation
silencing of mRNA), we apply local cross-attention inside each gene. The input
sequence for gene $i$ prepends a learnable token $\mathbf{z}_{cls}$:
\begin{equation}
    \mathbf{X}_i = \begin{bmatrix} \mathbf{z}_{cls} \\ \mathbf{z}_{m,i} \\
    \mathbf{z}_{c,i} \\ \mathbf{z}_{t,i} \end{bmatrix} \in \mathbb{R}^{4 \times d}.
\end{equation}

\vspace{0.15cm}\noindent\textbf{Unimodal Dropout.}
Masking a single modality enforces cross-modal redundancy. With probability
$p_{drop}$ at most one channel is masked:
\begin{equation}
    \tilde{\mathbf{z}}_{k,i} = \begin{cases}
      \text{0 now then }\mathbf{z}_{mask}^{(k)}, & \text{prob. } p_{drop} \\
      \mathbf{z}_{k,i}, & \text{otherwise}
   \end{cases} \quad \forall k \in \{m, c, t\},
\end{equation}
where $\mathbf{z}_{mask}^{(k)} \in \mathbb{R}^d$ is a learnable per-modality mask
token; capping at one masked channel lets the model impute the missing modality
from the remaining two. Because RNA-seq is the most information-dense modality and
tends to dominate, a higher dropout rate is applied to it specifically to prevent
the model from collapsing onto transcriptomics. \todo{Change 0 to
 $\mathbf{z}_{mask}^{(k)}$ when the masking token will be implemented and tested}

\vspace{0.15cm}\noindent\textbf{Intra-Gene Attention.}
Multi-head self-attention with a Pre-LayerNorm formulation,
\begin{equation}
    \mathbf{X}_i' = \mathbf{X}_i + \text{MHA}(\text{LN}(\mathbf{X}_i)),
\end{equation}
produces the unified gene representation $\mathbf{h}_i = \mathbf{X}_{i,0}'' \in
\mathbb{R}^d$ from the updated classification token.

\subsection{Graph Positional Encoding}
Beyond the attention bias below, each gene token is endowed with a graph-structural
positional encoding computed from the induced PPI subgraph of the selected genes.
We use Random-Walk Positional Encodings (RWPE)~\cite{dwivedi2022rwpe}: with
$\mathbf{A}$ the (binary, symmetric) adjacency and $\mathbf{P}=\mathbf{D}^{-1}
\mathbf{A}$ the random-walk transition matrix, the $k$-step return probabilities
\begin{equation}
    \mathbf{p}_i = \big[\, [\mathbf{P}^1]_{ii},\, [\mathbf{P}^2]_{ii},\, \dots,\,
    [\mathbf{P}^{d_{pe}}]_{ii} \,\big] \in \mathbb{R}^{d_{pe}}
\end{equation}
encode each gene's local topological role. The encoding is exposed as a modular
component (RWPE or Laplacian eigenvectors) so the positional scheme can itself be
tuned. 

\subsection{Global Transformer and Biased Attention}
\label{sec:globalattn}
To respect canonical pathways, we bias the global attention matrix with the
Shortest Path Distance (SPD) of the PPI network. Each head $h$ maps integer
distances to a learnable bias, bounded by $\tanh$ to preserve gradient stability:
\begin{equation}
    \mathbf{B}_{spatial}^{(h)} = \tanh\!\left(\text{Embed}^{(h)}(\mathbf{D}_{SPD})\right).
\end{equation}
A \texttt{[TUMOR\_CLS]} token aggregates information across all genes,
$\mathbf{H}^{(0)} = [\mathbf{t}_{cls}, \mathbf{h}_1, \dots, \mathbf{h}_N]$. In the
current model the structural bias is injected through separate row-wise softmax
terms (a dual-softmax formulation),
\begin{equation}
    \text{Attn}^{(h)} = \text{softmax}\!\left(\tfrac{\mathbf{Q}\mathbf{K}^\top}
    {\sqrt{d_k}}\right)\mathbf{V} + \text{softmax}\!\left(\mathbf{B}_{spatial}^{(h)}
    \right)\mathbf{V},
\end{equation}
following the row-wise scaling of Graphormer-Boosted~\cite{fadlallah2025}. A
canonical Graphormer formulation that places the bias \emph{inside} a single
softmax is a planned refinement (Sec.~\ref{sec:future}).

\subsection{Classification and Dampened Focal Loss}
The $K=8$ pathway tokens of the final layer are aggregated into a systemic tumor
state $\mathbf{u}_{tumor} = \text{MLP}_{agg}([\mathbf{t}_{p_1}^{(L)}, \dots,
\mathbf{t}_{p_K}^{(L)}]) \in \mathbb{R}^d$, and subtype logits are produced by a
bottleneck head ($d \to d/2$):
\begin{equation}
    \hat{\mathbf{y}} = \mathbf{W}_2\,\text{GELU}\!\left(\text{BatchNorm}
    (\mathbf{W}_1 \mathbf{u}_{tumor} + \mathbf{b}_1)\right).
\end{equation}
Optimization uses a Multi-Class Focal Loss~\cite{lin2017focal} with
$\gamma=2.0$, with square-root-dampened inverse class frequencies
$\alpha_c = \sqrt{1/N_c}$ to counter the extreme Normal-like imbalance without
inducing degenerate ``safe-guess'' behavior.

% =====================================================================
% INTERPRETABILITY
% =====================================================================
\section{Native Interpretability}
The architecture supports interpretation without external attribution methods:
\begin{itemize}
    \item \textbf{Global driver extraction:} the final-layer attention from the
    pathway tokens identifies the genes most heavily attended.
    \item \textbf{Deep drill-down:} for each driver gene, its $4\times4$
    intra-gene attention matrix isolates which omics modality drove the anomalous
    token state.
\end{itemize}

% =====================================================================
% EVALUATION METHODOLOGY
% =====================================================================
\section{Evaluation Methodology}
\label{sec:eval}
Because classical early-fusion models are strong in the $N \ll p$ regime, any claim
of deep-model benefit must survive a rigorous, identical protocol applied to all
methods. We therefore evaluate every model --- baselines and MOGFormer --- through
a single shared harness.

\vspace{0.15cm}\noindent\textbf{Cross-validation.}
We use repeated stratified $k$-fold cross-validation at the patient level: $5$
folds repeated $5$ times for $25$ estimates, stratified on PAM50 label. Five folds
(rather than ten) keep $\sim$7 Normal-like patients per held-out fold. Fold
indices are persisted so every model is scored on identical partitions, enabling
\emph{paired} statistical tests.

\vspace{0.15cm}\noindent\textbf{Leakage controls.}
Within each outer fold, HVG selection, imputation, normalization, \emph{and} the
induced PPI subgraph (hence the RWPE and SPD matrices) are recomputed on the
training patients only; no held-out patient influences any fitted statistic or the
graph. Hyperparameter tuning and early stopping use an inner stratified $20\%$
split of the training fold; the outer fold is never seen during model selection.

\vspace{0.15cm}\noindent\textbf{Statistics.}
We report macro-F1 as the primary metric (mean $\pm$ 95\% CI over the 25
estimates), together with per-class F1, balanced accuracy, Matthews correlation,
and per-class ROC/PR-AUC. Because cross-validation folds overlap, we additionally
report the Nadeau--Bengio corrected-resampled variance~\cite{nadeau2003} and, for
model-versus-model claims, the paired Wilcoxon signed-rank test on the 25 paired
per-fold scores, summarized through critical-difference diagrams~\cite{demsar2006}.

\vspace{0.15cm}\noindent\textbf{Baseline suite.}
All baselines run through the identical pipeline and folds: a stratified
\textit{Dummy} floor; a PAM50 nearest-centroid classifier on RNA (the
label-source reference~\cite{parker2009pam50}); penalized multinomial logistic
regression; a linear SVM; random forests; gradient-boosted trees (XGBoost); and
TabPFN-v2~\cite{hollmann2023tabpfn} on the top-$K$ genes. Each classical model is
evaluated under both early fusion (concatenated standardized features) and late
fusion (per-modality models with averaged class probabilities), and MOGONET is
retained as a direct GNN comparator~\cite{wang2021mogonet}.

% =====================================================================
% RESULTS
% =====================================================================
\section{Results}
\label{sec:results}
% NOTE: baseline numbers below are the current Phase-0 results; please confirm
% they are final before submission. The MOGFormer cross-validated result is still
% being computed and is intentionally left as "under evaluation".

\begin{table}[t]
    \centering
    \small
    \caption{PAM50 subtyping under the shared 25-estimate protocol (early fusion,
    all three modalities). Baseline figures are current Phase-0 results;
    MOGFormer is under evaluation through the identical harness.}
    \label{tab:baselines}
    \begin{tabular}{lc}
        \toprule
        Model & Macro-F1 \\
        \midrule
        Dummy (stratified) & \todo{floor} \\
        PAM50 nearest-centroid (RNA) & 0.738 \\
        Gradient-boosted trees (XGBoost) & \textbf{0.852} \\
        \midrule
        MOGFormer (ours) & \textit{under evaluation} \\
        \bottomrule
    \end{tabular}
\end{table}

\noindent\textbf{Classical early-fusion is a demanding baseline.} A well-tuned
XGBoost reaches a macro-F1 of $0.852$, well above the PAM50 nearest-centroid
reference ($0.738$); since PAM50 labels are themselves RNA-derived, the centroid
classifier is the most conceptually important reference point
(Table~\ref{tab:baselines}). This sets a stringent target for any deep model.

\vspace{0.15cm}\noindent\textbf{Multi-omics versus single-omics.} Across the
modality-contribution grid, RNA-only performance ($\sim$0.86 macro-F1) is
statistically indistinguishable from the full three-modality configuration, with
CNV and methylation adding little \emph{accuracy} over RNA alone
(Fig.~\ref{fig:results}). This motivates MOGFormer's design emphasis on
\emph{interpretable} cross-modal routing rather than accuracy gains from fusion
per se, and frames cross-modal contribution as a question best answered by
attention attribution rather than by held-out F1.

\vspace{0.15cm}\noindent\textbf{Which classes drive the error.} Per-class F1 for
the best classical model is high for Basal ($0.98$), Luminal A ($0.94$), HER2
($0.89$), and Luminal B ($0.88$), but collapses for Normal-like ($0.51$). The
dominant confusions are Luminal A $\leftrightarrow$ Luminal B and Normal-like
$\to$ Luminal A. Normal-like ($n=35$) is therefore the universal bottleneck, and
we report per-class F1 at every step to prevent the macro average from masking
this collapse.

\vspace{0.15cm}\noindent\textbf{Feature-count sensitivity.} Performance is stable
for $N \in [100, 500]$ selected genes, defining the operating regime; we fix
$N=250$ to balance this plateau against the $O(N^2)$ cost of global attention.

\begin{figure}[t]
    \centering
    % \includegraphics[width=0.95\linewidth]{figures/C00_baseline_vs_references.pdf}
    \figplaceholder{Figure 2 placeholder --- MOGFormer macro-F1 ($\pm$95\% CI)
    against reference lines (XGBoost 0.852, PAM50-centroid 0.738), plus the
    per-class F1 panel emphasizing Normal-like.}
    \caption{MOGFormer versus reference models under the shared protocol, with
    per-class breakdown. \todo{Insert C00/F01 once the 25-estimate MOGFormer run
    completes.}}
    \label{fig:results}
\end{figure}

% =====================================================================
% FUTURE WORK
% =====================================================================
\section{Ongoing Refinements and Future Work}
\label{sec:future}

\subsection{Targeted Architectural Refinements}
Guided by the gap to the classical baseline, four single-variable refinements are
under controlled evaluation, each applied and measured in isolation through the
shared harness: (i) a per-gene identity embedding $\mathbf{G}_i$ added to every
modality token, so the value path can distinguish genes of equal scaled magnitude;
(ii) a canonical Graphormer attention that places the structural bias inside a
single softmax rather than the dual-softmax of Sec.~\ref{sec:globalattn}; (iii)
periodic (PLR) numerical embeddings~\cite{gorishniy2022} replacing the per-modality
MLP tokenizers; and (iv) the learnable per-modality mask token for unimodal
dropout. \todo{Report deltas once each single-change run completes; these are
hypotheses, not yet results.}

\subsection{Self-Supervised Pan-Cancer Pre-Training}
To mitigate the $N \ll p$ constraint, we plan a self-supervised pre-training stage
on unlabeled pan-cancer TCGA cohorts, building an oncogenic foundation
representation before downstream adaptation. A random $15\%$ subset $\mathcal{M}$
of genes has its three modalities replaced by a learnable mask token,
$\mathbf{h}_i^{(0)} = \mathbf{z}_{mask}\ \forall i \in \mathcal{M}$; the masked
sequence is propagated through the structurally biased layers and a regression head
reconstructs the original normalized values, optimized by a masked MSE,
\begin{equation}
    \mathcal{L}_{MGM} = \frac{1}{|\mathcal{M}|} \sum_{i \in \mathcal{M}}
    \sum_{k \in \{m,c,t\}} \left( x_{k,i} - \hat{x}_{k,i} \right)^2.
\end{equation}
After convergence, the reconstruction head is discarded, the tokenizers,
intra-gene mini-transformers, and lower global layers are frozen, and the pathway
tokens plus top classification layers are fine-tuned with the dampened focal loss
--- reducing the sample complexity required on rare subtypes. \todo{This section
describes planned work; mark clearly as such or omit for a results-focused
submission.}

% =====================================================================
% METHODS
% =====================================================================
\section*{Methods}

\subsection*{Graph Construction and Spatial Bias Embedding}
The biological prior is the Human PPI network from STRING~\cite{szklarczyk2023string}.
Let $G=(V,E)$ with $V$ the per-fold selected genes ($N$) and $E$ the interactions
with combined score $\geq 400$, mapped from STRING protein identifiers to HGNC
symbols via the alias table. The graph is induced on the selected genes and rebuilt
per fold. The Shortest Path Distance matrix $\mathbf{D}_{SPD} \in \mathbb{R}^{N
\times N}$ is computed with a cutoff $D_{max}$; unreachable pairs are assigned
$D_{max}+1$. Each integer distance maps to a learnable per-head scalar bias bounded
by $\tanh$.

\subsection*{Model Optimization and Hyperparameters}
MOGFormer is trained with AdamW ($\eta = 10^{-4}$, weight decay $10^{-4}$), token
dimension $d=128$, $K=8$ latent pathway tokens, and the Multi-Class Focal Loss
($\gamma=2.0$, square-root-dampened class weights). Early stopping monitors inner-
validation macro-F1. \todo{Align the following to the code you are running: the
current configuration uses $L=1$ global transformer layer and $D_{max}=10$ }

% =====================================================================
% REFERENCES
% =====================================================================
% NOTE: bibliographic details below are landmark works cited throughout; please
% verify exact volume/page/DOI against your reference manager before submission.
\normalsize
\begin{thebibliography}{99}
\small
\bibitem{parker2009pam50} Parker, J.S., \textit{et al.} Supervised risk predictor
of breast cancer based on intrinsic subtypes. \textit{J. Clin. Oncol.} (2009).
\bibitem{tcga2012brca} The Cancer Genome Atlas Network. Comprehensive molecular
portraits of human breast tumours. \textit{Nature} (2012).
\bibitem{ying2021graphormer} Ying, C., \textit{et al.} Do Transformers Really
Perform Bad for Graph Representation? \textit{NeurIPS} (2021).
\bibitem{fadlallah2025} Fadlallah, S.A. Graphormer-Boosted: Molecular Property
Prediction with Enhanced Graph Spatial and Edge Encodings (2025).
\bibitem{dwivedi2022rwpe} Dwivedi, V.P., \textit{et al.} Graph Neural Networks with
Learnable Structural and Positional Representations. \textit{ICLR} (2022).
\bibitem{gorishniy2022} Gorishniy, Y., Rubachev, I., Babenko, A. On Embeddings for
Numerical Features in Tabular Deep Learning. \textit{NeurIPS} (2022).
\bibitem{lin2017focal} Lin, T.-Y., \textit{et al.} Focal Loss for Dense Object
Detection. \textit{ICCV} (2017).
\bibitem{wang2021mogonet} Wang, T., \textit{et al.} MOGONET integrates multi-omics
data using graph convolutional networks. \textit{Nat. Commun.} (2021).
\bibitem{hollmann2023tabpfn} Hollmann, N., \textit{et al.} TabPFN: A Transformer
That Solves Small Tabular Classification Problems in a Second. \textit{ICLR} (2023).
\bibitem{szklarczyk2023string} Szklarczyk, D., \textit{et al.} The STRING database:
protein--protein association networks. \textit{Nucleic Acids Res.} (2023).
\bibitem{nadeau2003} Nadeau, C., Bengio, Y. Inference for the Generalization Error.
\textit{Machine Learning} (2003).
\bibitem{demsar2006} Dem\v{s}ar, J. Statistical Comparisons of Classifiers over
Multiple Data Sets. \textit{J. Mach. Learn. Res.} (2006).
\end{thebibliography}

\end{document}